\section{Scene Classification using Convolutional Neural Networks}

The usage of Deep Convolutional Neural Networks (CNNs) is currently the way to go to solve image recognition and classification tasks. [https://papers.nips.cc/paper/4824-imagenet-classification-with-deep-convolutional-neural-networks.pdf]. CNNs, inspired by biological processes, perform so well because the whole processing, assuming that the image is of the correct size, is done by a single neural network with multiple layers of neurons and millions of learned parameters. In case of a classification task the output of the CNN is a probability distribution over the classes or at least a probability-like confidence value for each class.

Beside choosing a good network architecture and training algorithm the training data is very important. The Place205 dataset contains 2.5 million images classified into 205 different scene categories. A computer vision challenge using this dataset resulted in multiple CNNs available to be used in scene classification tasks.